% ID: SEL-D0049-OB-E01
% ESTADO: cotejada. Segunda pasada visual y borrador compilado revisados.
% Fuente: Selectividad/sel_2013_matematicas_aplicadas/junio/reserva b/EXAMEN 1-ANDALUCÍA.pdf
% SHA-256: ec6fd00d0daee547aa6128f3f6e1099278ee551f7f80917b331a1772764317b3
% Página: 2; opción B; ejercicio original 1.

\begin{enumerate}[label=\alph*)]
\item \textbf{(1 punto)} En un problema de programación lineal, la región factible es la región acotada cuyos vértices son $A(2,-1)$, $B(-1,2)$, $C(1,4)$ y $D(5,0)$. La función objetivo es la función $f(x,y)=2x+3y+k$, cuyo valor máximo, en dicha región, es igual a 19. Calcule el valor de $k$ e indique dónde se alcanza el máximo y dónde el mínimo.
\item \textbf{(1.5 puntos)} Sean las matrices
\[
A=\begin{pmatrix}1&-2&3\end{pmatrix},\quad
B=\begin{pmatrix}2\\-1\\1\end{pmatrix},\quad
C=\begin{pmatrix}2&0&-1\\1&1&-1\\1&3&2\end{pmatrix}.
\]
Resuelva, si es posible, la ecuación matricial $B\cdot A+2X=C$.
\end{enumerate}
